\section{Theoretical Background}
\label{sec:theoretical_background}

\subsection{Optimization Theory}
\label{subsec:optimization_theory}

\subsubsection{Black-Box Optimization}
\label{subsubsec:black_box_optimization}

A black-box optimization problem is one in which the objective function has no known analytical form and no gradient information is available to the
optimizer\cite{Audet2017Derivative-Free}. The function can only be queried at candidate input points, returning a scalar or vector output. This
setting arises naturally in engineering design when objectives are evaluated through computationally expensive procedures, making gradient computation
either impossible or impractical. Because each function evaluation may be costly, black-box optimizers must find high-quality solutions within a
limited evaluation budget, placing emphasis on sample efficiency\cite{Jones2022Constrained}. The gripper configuration problems addressed in this
thesis fall squarely within this category: objective values are computed through a series of geometric and numerical calculations applied to the
candidate gripper configuration, and no closed-form expression exists that directly maps the decision variables to either the contact quality score or
the grasp diversity metric.

\subsubsection{Multi-Objective Optimization}
\label{subsubsec:multi_objective_optimization}

Multi-objective optimization (MOO) refers to the simultaneous optimization of two or more conflicting objective functions over a shared decision
space\cite{Deb2011Multi-objective}. Formally, a multi-objective minimization problem is defined as:

\begin{equation}
    \min_{\mathbf{x} \in \mathcal{X}} \mathbf{f}(\mathbf{x}) =
    \bigl(f_1(\mathbf{x}),\, f_2(\mathbf{x}),\, \ldots,\, f_k(\mathbf{x})\bigr)
\end{equation}

where $\mathbf{x}$ is the decision variable vector, $\mathcal{X}$ is the feasible decision space, and $k \geq 2$ is the number of objectives. When
objectives conflict, no single solution minimizes all objectives simultaneously. Instead, a set of trade-off solutions exists, collectively forming
the Pareto front. The goal of a multi-objective optimizer is therefore not to find a single best solution but to approximate this trade-off surface as
accurately and completely as possible\cite{Deb2011Multi-objective,Gunantara2018A}.

\subsubsection{Pareto Optimality}
\label{subsubsec:pareto_optimality}

The concept of Pareto optimality provides the theoretical foundation for comparing solutions in a multi-objective setting. A solution
$\mathbf{x}^{(a)}$ is said to \emph{dominate} $\mathbf{x}^{(b)}$, written $\mathbf{x}^{(a)} \prec \mathbf{x}^{(b)}$, if and only if:

\begin{equation}
    f_i(\mathbf{x}^{(a)}) \leq f_i(\mathbf{x}^{(b)}) \quad \forall\, i \in
    \{1,\ldots,k\}
    \quad \text{and} \quad
    f_j(\mathbf{x}^{(a)}) < f_j(\mathbf{x}^{(b)}) \quad \text{for at least
    one } j
\end{equation}

A solution is called \emph{Pareto optimal} if no other feasible solution dominates it. The set of all Pareto-optimal solutions in the decision space
is the \emph{Pareto set}, and its image in the objective space is the \emph{Pareto front}. In practice, algorithms compute an approximation of the
Pareto front, and the quality of this approximation is assessed using performance indicators such as the hypervolume
indicator\cite{Audet2018Performance,Shang2021A}.

\subsubsection{Evolutionary Algorithms}
\label{subsubsec:evolutionary_algorithms}

Evolutionary algorithms (EAs) are a family of population-based metaheuristics inspired by the principles of biological evolution. A population of
candidate solutions is maintained across discrete generations. At each generation, new candidate solutions are generated through variation operators
typically crossover, which recombines pairs of existing solutions, and mutation, which introduces random perturbations and a selection mechanism
retains the most promising individuals for the next generation based on their objective values. This loop continues until a termination criterion is
met, such as a maximum number of function evaluations or convergence of the population\cite{Abulail2025Enhanced}.

EAs are well suited to black-box multi-objective optimization because they require no gradient information, naturally maintain a diverse population of
trade-off solutions, and are robust to noisy or discontinuous objective landscapes. Multi-objective evolutionary algorithms (MOEAs) extend the basic
EA framework with selection mechanisms based on Pareto dominance or decomposition, enabling them to approximate the Pareto front in a single
run\cite{Ishibuchi2017Performance,Li2015An}.

\subsubsection{Bayesian Optimization}
\label{subsubsec:bayesian_optimization}

Bayesian optimization (BO) is a sequential model-based strategy for the global optimization of expensive black-box functions. A probabilistic
surrogate model most commonly a Gaussian process (GP) is fitted to the observations collected so far\cite{Frazier2018A}. The GP provides both a mean
prediction and an uncertainty estimate at unobserved locations, which together define an \emph{acquisition function} that balances exploration of
uncertain regions with exploitation of promising ones\cite{Frazier2018A,Candelieri2021A}. The point that maximizes the acquisition function is
selected as the next evaluation, after which the surrogate is updated and the process repeats\cite{Wang2022Recent}.

\subsubsection{Hyperparameter Tuning}
\label{subsubsec:hyperparameter_tuning}

Optimization algorithms expose a set of hyperparameters such as population size, crossover probability, and mutation rate that control their search
behaviour but are not adapted during the optimization run itself. The choice of hyperparameters can substantially affect algorithm performance, and
using default values may systematically favour or disadvantage certain algorithms in a comparative
study\cite{Sivaprasad2019Optimizer,Ishibuchi2022Difficulties}.

Hyperparameter tuning addresses this by searching the hyperparameter space for settings that improve algorithm performance on the problem at hand
before the main optimization run begins. In this thesis, algorithm hyperparameters are initialized from the default values provided by the
\texttt{pymoo} framework\cite{pymoo}. Where the optimization pipeline is configured to do so, a Bayesian optimization procedure is executed prior to
the main run to replace these defaults with problem-specific settings, reducing the risk of performance bias from incidental parameter
choices\cite{Probst2018Tunability:}.

%-------------------------------------------------------------------
\subsection{Multi-Objective Evolutionary Algorithms}
\label{subsec:moea}

The candidate algorithms evaluated in this thesis are all members of the multi-objective evolutionary algorithm (MOEA) family. Each follows the
general EA loop described in Section~\ref{subsubsec:evolutionary_algorithms} but differs in how it guides selection toward a well-distributed
approximation of the Pareto front. Brief introductions to each algorithm are given below.

\subsubsection{NSGA-II}
\label{subsubsec:nsga2}

The Non-dominated Sorting Genetic Algorithm II (NSGA-II) is among the most widely adopted MOEAs. Its selection mechanism is based on two criteria
applied in sequence. First, the entire population is sorted into non-domination ranks: all non-dominated solutions receive rank 1, they are removed,
and the process repeats until every individual has a rank. Second, within the same rank, solutions are ordered by \emph{crowding distance}, a density
estimate that measures how isolated a solution is from its neighbours in the objective space. Solutions with a higher crowding distance are preferred,
encouraging spread along the Pareto front. NSGA-II uses an elitist strategy in which the combined parent and offspring population of size $2N$ is
truncated to $N$ using these two criteria. Its computational complexity per generation is $\mathcal{O}(kN^2)$, where $k$ is the number of objectives
and $N$ is the population size\cite{nsga2}.

\subsubsection{NSGA-III}
\label{subsubsec:nsga3}

NSGA-III extends NSGA-II to problems with three or more objectives (many-objective problems), where crowding distance becomes an unreliable diversity
measure. Instead of crowding distance, NSGA-III uses a set of structured \emph{reference points} distributed uniformly on a normalized hyperplane in
the objective space. Solutions are associated with their closest reference point, and selection preference is given to solutions that are associated
with under-represented reference points and lie closest to their reference line. This mechanism explicitly promotes diversity across the objective
space and has been shown to outperform NSGA-II on problems with three or more objectives\cite{nsga3}.

\subsubsection{MOEA/D}
\label{subsubsec:moead}

MOEA/D (Multi-Objective Evolutionary Algorithm based on Decomposition) takes a fundamentally different approach from Pareto-ranking methods. The
multi-objective problem is decomposed into $N$ scalar subproblems using a set of evenly distributed weight vectors and an aggregation function such as
the weighted Chebyshev scalarization. Each subproblem is optimized simultaneously, with solutions sharing information within a defined neighbourhood
of weight vectors. This neighbourhood-based collaborative search makes MOEA/D computationally efficient and well suited to problems where the Pareto
front has a regular structure. Its complexity per generation is lower than NSGA-II, as it avoids global non-dominated sorting\cite{pymoo}.

\subsubsection{SMS-EMOA}
\label{subsubsec:smsemoa}

The S-Metric Selection Evolutionary Multi-Objective Algorithm (SMS-EMOA) uses the hypervolume indicator directly as its selection criterion. It
operates as a steady-state algorithm: at each step, a single offspring is generated, added to the population, and the individual that contributes the
least to the population hypervolume is removed. This ensures that the population always converges toward solutions that maximize the dominated
hypervolume with respect to a reference point. SMS-EMOA is particularly well suited to problems where solution quality is to be measured by the
hypervolume indicator, as its selection pressure is directly aligned with that metric\cite{smsemoa}.

\subsubsection{CMOPSO}
\label{subsubsec:cmopso}

The Competitive Multi-Objective Particle Swarm Optimizer (CMOPSO) adapts the particle swarm optimization (PSO) framework to the multi-objective
setting. In standard PSO, each particle moves through the decision space guided by its personal best position and a global best position. CMOPSO
introduces a competitive mechanism in which particles compete in pairs within the swarm; the losing particle updates its velocity by learning from the
winning particle rather than from a fixed global best. An external archive of non-dominated solutions is maintained to provide guidance and preserve
diversity across the Pareto front. This competitive selection strategy avoids premature convergence to a single region of the Pareto front and has
been shown to achieve fast convergence while remaining competitive with dominance-based MOEAs on standard benchmark problems\cite{cmopso}.

%-------------------------------------------------------------------
\subsection{Gripper Theory and Geometric Representation}
\label{subsec:gripper_theory}

\subsubsection{Grasp Theory}
\label{subsubsec:grasp_theory}

A grasp is a configuration of contact points between a gripper and a workpiece such that the gripper can resist external forces and torques applied to
the object\cite{Prattichizzo2014Robot}. The classical criterion for grasp quality is \emph{force closure}: a grasp is force-closed if the contact
forces available at the contact points can produce any net wrench on the object, meaning that arbitrary disturbances can be
resisted\cite{Nguyen1986Constructing}. Force closure requires a sufficient number of contacts with appropriate friction properties and geometric
placement.

\subsubsection{Suction Cup Gripper and Grasp Diversity}
\label{subsubsec:suction_cup_theory}

The \emph{diversity} of a grasp set can be quantified by combining two complementary measures of spatial spread.

\paragraph{Geodesic Distance as a Diversity Measure.}
The geodesic distance between two points on a surface is the length of the shortest path between them that lies entirely on the surface. As a
diversity measure, geodesic distance captures surface-level spread that Euclidean distance in ambient space does not reflect for curved geometries. A
larger mean geodesic distance between grasp contact points indicates broader sampling of the workpiece surface, which is desirable for robustness to
workpiece pose variation\cite{Peyre2010Geodesic}.

\paragraph{Spatial Diversity via Mean Pairwise Euclidean Distance.}
The spatial spread of a grasp set can be quantified by the mean pairwise Euclidean distance between grasp positions in three-dimensional space,
normalized by a characteristic object scale $s$ to ensure comparability across objects of different sizes. The characteristic scale is defined as the
Euclidean norm of the axis-aligned bounding-box extents:

\begin{equation}
    s = \|\mathbf{e}\|_2
\end{equation}

where $\mathbf{e} \in \mathbb{R}^3$ contains the bounding-box side lengths. The normalized spatial diversity is then:

\begin{equation}
    D_{\mathrm{pos}}(\mathbf{p}) =
    \frac{1}{s}
    \cdot
    \frac{1}{n(n-1)}
    \sum_{\substack{i,j=1 \\ i \neq j}}^{n}
    \|\mathbf{y}_i - \mathbf{y}_j\|_2
\end{equation}

where $\mathbf{y}_i \in \mathbb{R}^3$ is the position of the $i$-th grasp contact point and $n$ is the total number of grasps considered. Normalizing
by $s$ ensures that the metric is scale-invariant, so that diversity values remain comparable across workpieces of differing sizes. This measure
complements the geodesic measure by reflecting absolute spatial separation rather than surface-constrained distance.

\subsubsection{Two-Finger Gripper and Contour Representation}
\label{subsubsec:two_finger_theory}

The two-finger gripper applies gripping forces through two opposing jaw surfaces whose contour geometry determines how well the gripper conforms to
the workpiece. Designing the jaw contour requires a representation that is compact enough to serve as a low-dimensional decision space for
optimization, yet expressive enough to capture smooth, physically realizable shapes. Spline parameterization is a well-established approach for this
purpose.

\paragraph{B-Splines.}
A B-spline curve of degree $p$ is defined by a set of $n+1$ control points $\{\mathbf{P}_0, \mathbf{P}_1, \ldots, \mathbf{P}_n\}$ and a non-decreasing
knot vector $\mathbf{t} = \{t_0, t_1, \ldots, t_{n+p+1}\}$ \cite{Eilers1996Flexible}. The curve is expressed as:

\begin{equation}
    \mathbf{C}(u) = \sum_{i=0}^{n} N_{i,p}(u)\, \mathbf{P}_i,
    \quad u \in [t_p,\, t_{n+1}]
\end{equation}

where $N_{i,p}(u)$ are the B-spline basis functions of degree $p$, defined recursively by the Cox--de Boor recurrence
relation\cite{Eilers1996Flexible}. B-spline curves possess the \emph{local support} property: moving a single control point affects the curve only in
its neighbourhood, making them well suited to constrained shape optimization\cite{Hasan2024B-spline,Eilers1996Flexible}. They also guarantee $C^{p-1}$
continuity everywhere except at repeated knots, ensuring smooth contour shapes without explicit smoothness constraints in the optimizer.

\paragraph{Spline-Based Shape Parameterization.}
Spline-based shape parameterization represents a geometric contour through a small set of control points, reducing an otherwise infinite-dimensional
shape space to a finite-dimensional parameter vector. Boundary control points are typically fixed to satisfy geometric constraints, while interior
control points are free to vary and serve as the design variables. The number of free control points governs the trade-off between geometric
expressiveness and search space dimensionality: more control points allow more complex shapes but increase the dimensionality of the optimization
problem\cite{Hasan2024B-spline,Chapelier2022Spline-based}.